El proyecto busca predecir de forma temprana el retiro de estudiantes en entornos de aprendizaje virtual mediante datos académicos y de participación. Esta versión presenta la revisión y preparación de OULAD (\textit{Open University Learning Analytics Dataset}), siguiendo CRISP-DM, una metodología que organiza el proceso de minería de datos. Se examinaron siete tablas que incluyen 32\,593 inscripciones de 28\,785 personas, y se construyeron indicadores para los días 7, 14, 28, 42 y 56 del curso.

En cada corte se utilizó información registrada hasta ese momento y se excluyeron los retiros ya ocurridos. El evento que se busca predecir es el retiro administrativo posterior al corte y hasta el final de la presentación; la inactividad se conserva como indicador complementario. Las inscripciones elegibles pasaron de 29\,076 en el primer corte a 26\,506 en el último. Se observaron diferencias en la actividad, las entregas y su relación con el retiro. Estas comparaciones deben leerse teniendo en cuenta que cambian tanto el grupo analizado como el tiempo restante de seguimiento. En el corte 7 no había evaluaciones con vencimiento dentro de la ventana, aunque ya existían entregas; por ello, no entregar todavía no equivale necesariamente a incumplir una actividad.

Los análisis se ejecutaron desde las fuentes y se acompañaron de comprobaciones independientes de clics, días activos, variabilidad entre días activos y entregas. Los resultados dejan una base documentada para comparar modelos, pero aún no permiten elegir una ventana óptima ni afirmar capacidad predictiva o efectividad educativa. El entrenamiento, la interpretación mediante SHAP (atribuciones de las variables a las predicciones del modelo) y el prototipo de visualización permanecen pendientes.

\textbf{Palabras clave}: abandono estudiantil, OULAD, analítica de aprendizaje, ventanas temporales, retiro administrativo.
